% from ega3-7-fr.tex
\oldpage[III]{220}
en supposant seulement que $(X,\mathscr O_X)$ soit un espace annelé en anneaux
locaux, $\mathscr J$ un Idéal de $\mathscr O_X$ tel que $\mathscr J^n=0$, que
l'espace annelé $X_0=(X,\mathscr O_X/\mathscr J)$ soit un schéma affine et enfin
que les Idéaux $\mathscr J^k/\mathscr J^{k+1}$ soient des
$\mathscr O_{X_0}$-Modules quasi-cohérents. Il suffit en effet d'utiliser
(1.8.1) (cf. $(\mathbf{Err}_{\mathrm{II}})$) au lieu de (2.2.4) dans le
raisonnement.

\noindent\textbf{$(\mathbf{Err}_{\mathrm{III}},\,9)$} Dans (I, 5.3.1), ligne
11 de la p. 132, insérer, après « ou $\Delta_X$ », « ou $\Delta_\varphi$ si
$\varphi:X\to S$ est le morphisme structural ».

\noindent\textbf{$(\mathbf{Err}_{\mathrm{III}},\,10)$} Dans (I, 5.3.9), la
démonstration est insuffisante, car elle ne prouve pas que $\Delta_X(X)$ soit
localement fermé dans $X\times_SX$. Pour avoir une démonstration correcte, il
suffit d'utiliser (4.2.4, $a$)) : pour tout $x\in X$ et tout voisinage affine
$U$ de $x$ dans $X$, $U\times_SU$ est un voisinage affine de $\Delta_X(x)$ ;
compte tenu de (5.3.16) (dont la démonstration n'utilise que la définition
(5.3.1.1)), on est ramené à démontrer (5.3.9) lorsque
$S=\operatorname{Spec}(B)$ et $X=\operatorname{Spec}(A)$ sont des schémas
affines ; il est clair alors en vertu de (5.3.1.1) que $\Delta_X$ correspond à
l'homomorphisme canonique $A\otimes_BA\to A$ qui transforme $x\otimes y$ en
$xy$ ; cet homomorphisme étant surjectif, $\Delta_X$ est dans ce cas une
immersion fermée (4.2.3), ce qui achève la démonstration.

\noindent\textbf{$(\mathbf{Err}_{\mathrm{III}},\,11)$} Dans (I, 7.1.15) et
(I, 7.1.16), lignes 7 et 15 du bas de la p. 158, supprimer
« \emph{géométriques} ».

\noindent\textbf{$(\mathbf{Err}_{\mathrm{III}},\,12)$} Remplacer (I, 7.4.7)
par : On étend (par abus de langage) les définitions de (7.4.1) au cas où
$X$ est un préschéma \emph{réduit} dont tout point admet un voisinage ouvert
n'ayant qu'un nombre \emph{fini} de composantes irréductibles ; il résulte
alors de (7.3.4) et (7.4.6) que, pour un $\mathscr O_X$-Module quasi-cohérent de
type fini $\mathscr F$, dire que $\mathscr F$ est un \emph{faisceau de torsion}
équivaut à dire que $\operatorname{Supp}(\mathscr F)$ \emph{ne contient aucune
composante irréductible de $X$}.

\noindent\textbf{$(\mathbf{Err}_{\mathrm{III}},\,13)$} Dans (I, 10.11.7),
lignes 17 à 23 de la p. 205, la fin de la démonstration depuis « Reste à
prouver\ldots{} » est inutile en vertu de (10.11.6), les faisceaux considérés
étant cohérents par $(0,\,5.3.5)$.

\noindent\textbf{$(\mathbf{Err}_{\mathrm{III}},\,14)$} Dans (II, 1.7.8),
après la ligne 10 du bas de la p. 16, ajouter : lorsque
$\mathscr E=\mathscr O_S^n$, on écrit aussi $\mathbf V_S^n$ au lieu de
$\mathbf V(\mathscr E)$ ; si en outre $S=\operatorname{Spec}(A)$ est affine,
on écrit $\mathbf V_A^n$ au lieu de $\mathbf V_S^n$ ; on a
$\mathbf V_A^n=\operatorname{Spec}(A[T_1,\ldots,T_n])$ où les $T_i$ sont des
indéterminées.

\noindent\textbf{$(\mathbf{Err}_{\mathrm{III}},\,15)$} Dans (II, 1.7.10),
ligne 10 de la p. 17, supprimer « \emph{géométriques} » ; lignes 12 et 16--17
de la p. 17, supprimer « géométrique ».

\noindent\textbf{$(\mathbf{Err}_{\mathrm{III}},\,16)$} Dans (II, 1.7.12),
ajouter : On pose de même :
\[
  \mathbf V(\mathscr O_S^n)=\mathbf V_S^n=S[T_1,\ldots,T_n].
\]

\noindent\textbf{$(\mathbf{Err}_{\mathrm{III}},\,17)$} Dans (II, 4.2.6),
ligne 7 de la p. 75, supprimer « \emph{géométriques} » ; ligne 8 de la p.
75, supprimer « \emph{géométrique} ».

\noindent\textbf{$(\mathbf{Err}_{\mathrm{III}},\,18)$} Dans (II, 4.6.13),
ligne 7 de la p. 92, le raisonnement doit être précisé, car avec les notations
de (4.4.10), il faut ici prouver que l'immersion $\Gamma_f$ est
\emph{quasi-compacte}, afin d'appliquer (i bis). En vertu de (4.6.4), pour
prouver que $\mathscr L$ est ample relativement à $f$, on peut se borner au
cas où $Y$ est affine. Notons d'autre part que dans chacune des hypothèses de
(v), $f$ est quasi-compact (I, 6.6.4). D'autre part, si $g$ est
